\documentclass[10pt,a4paper]{amsart}
\usepackage{amsmath,amssymb,amsthm,graphicx}
\usepackage{verbatim}
\newenvironment{reference}{\comment}{\endcomment}
%\newenvironment{reference}{}{}
\newenvironment{slogan}{\comment}{\endcomment}
\newenvironment{history}{\comment}{\endcomment}

% For commutative diagrams we use Xy-pic
\usepackage[all]{xy}

% We use 2cell for 2-commutative diagrams.
\xyoption{2cell}
% Theorem environments.
%
\theoremstyle{plain}
\newtheorem{theorem}[subsection]{Theorem}
\newtheorem{proposition}[subsection]{Proposition}
\newtheorem{lemma}[subsection]{Lemma}

\theoremstyle{definition}
\newtheorem{definition}[subsection]{Definition}
\newtheorem{example}[subsection]{Example}
\newtheorem{exercise}[subsection]{Exercise}
\newtheorem{situation}[subsection]{Situation}

\theoremstyle{remark}
\newtheorem{remark}[subsection]{Remark}
\newtheorem{remarks}[subsection]{Remarks}

\numberwithin{equation}{subsection}

% Macros
%
\def\lim{\mathop{\mathrm{lim}}\nolimits}
\def\colim{\mathop{\mathrm{colim}}\nolimits}
\def\Spec{\mathop{\mathrm{Spec}}}
\def\Hom{\mathop{\mathrm{Hom}}\nolimits}
\def\Ext{\mathop{\mathrm{Ext}}\nolimits}
\def\SheafHom{\mathop{\mathcal{H}\!\mathit{om}}\nolimits}
\def\SheafExt{\mathop{\mathcal{E}\!\mathit{xt}}\nolimits}
\def\Sch{\mathit{Sch}}
\def\Mor{\mathop{\mathrm{Mor}}\nolimits}
\def\Ob{\mathop{\mathrm{Ob}}\nolimits}
\def\Sh{\mathop{\mathit{Sh}}\nolimits}
\def\NL{\mathop{N\!L}\nolimits}
\def\CH{\mathop{\mathrm{CH}}\nolimits}
\def\proetale{{pro\text{-}\acute{e}tale}}
\def\etale{{\acute{e}tale}}
\def\QCoh{\mathit{QCoh}}
\def\Ker{\mathop{\mathrm{Ker}}}
\def\Im{\mathop{\mathrm{Im}}}
\def\Coker{\mathop{\mathrm{Coker}}}
\def\Coim{\mathop{\mathrm{Coim}}}

% Boxtimes
%
\DeclareMathSymbol{\boxtimes}{\mathbin}{AMSa}{"02}

%
% Macros for moduli stacks/spaces
%
\def\QCohstack{\mathcal{QC}\!\mathit{oh}}
\def\Cohstack{\mathcal{C}\!\mathit{oh}}
\def\Spacesstack{\mathcal{S}\!\mathit{paces}}
\def\Quotfunctor{\mathrm{Quot}}
\def\Hilbfunctor{\mathrm{Hilb}}
\def\Curvesstack{\mathcal{C}\!\mathit{urves}}
\def\Polarizedstack{\mathcal{P}\!\mathit{olarized}}
\def\Complexesstack{\mathcal{C}\!\mathit{omplexes}}
% \Pic is the operator that assigns to X its picard group, usage \Pic(X)
% \Picardstack_{X/B} denotes the Picard stack of X over B
% \Picardfunctor_{X/B} denotes the Picard functor of X over B
\def\Pic{\mathop{\mathrm{Pic}}\nolimits}
\def\Picardstack{\mathcal{P}\!\mathit{ic}}
\def\Picardfunctor{\mathrm{Pic}}
\def\Deformationcategory{\mathcal{D}\!\mathit{ef}}
\usepackage{cleveref}
\crefname{originalreference}{original reference}{original references}
\Crefname{originalreference}{Original reference}{Original references}
\makeatletter
\newlabel{topologies-lemma-sheaf-property-fpqc}{{022H}{source}{}{sourcepointer-1}{}}
\newlabel{topologies-lemma-sheaf-property-fpqc@cref}{{[originalreference][0][]022H}{[1][0][]source}{}{}{}}
\makeatother
\crefname{originalreference}{original reference}{original references}
\Crefname{originalreference}{Original reference}{Original references}
\makeatletter
\newlabel{spaces-lemma-properties-diagonal}{{02X4}{source}{}{sourcepointer-2}{}}
\newlabel{spaces-lemma-properties-diagonal@cref}{{[originalreference][0][]02X4}{[1][0][]source}{}{}{}}
\makeatother
\crefname{originalreference}{original reference}{original references}
\Crefname{originalreference}{Original reference}{Original references}
\makeatletter
\newlabel{descent-lemma-descending-property-isomorphism}{{02L4}{source}{}{sourcepointer-3}{}}
\newlabel{descent-lemma-descending-property-isomorphism@cref}{{[originalreference][0][]02L4}{[1][0][]source}{}{}{}}
\makeatother
\crefname{originallemma}{original lemma}{original lemmas}
\Crefname{originallemma}{Original lemma}{Original lemmas}
\makeatletter
\newlabel{lemma-topology-points}{{03BX}{source}{}{sourcepointer-4}{}}
\newlabel{lemma-topology-points@cref}{{[originallemma][0][]03BX}{[1][0][]source}{}{}{}}
\makeatother
\crefname{originallemma}{original lemma}{original lemmas}
\Crefname{originallemma}{Original lemma}{Original lemmas}
\makeatletter
\newlabel{lemma-cover-by-union-affines}{{03FX}{source}{}{sourcepointer-5}{}}
\newlabel{lemma-cover-by-union-affines@cref}{{[originallemma][0][]03FX}{[1][0][]source}{}{}{}}
\makeatother
\crefname{originallemma}{original lemma}{original lemmas}
\Crefname{originallemma}{Original lemma}{Original lemmas}
\makeatletter
\newlabel{lemma-separated-cover}{{03FZ}{source}{}{sourcepointer-6}{}}
\newlabel{lemma-separated-cover@cref}{{[originallemma][0][]03FZ}{[1][0][]source}{}{}{}}
\makeatother
\crefname{originallemma}{original lemma}{original lemmas}
\Crefname{originallemma}{Original lemma}{Original lemmas}
\makeatletter
\newlabel{lemma-points-cartesian}{{03H4}{source}{}{sourcepointer-7}{}}
\newlabel{lemma-points-cartesian@cref}{{[originallemma][0][]03H4}{[1][0][]source}{}{}{}}
\makeatother
\crefname{originalreference}{original reference}{original references}
\Crefname{originalreference}{Original reference}{Original references}
\makeatletter
\newlabel{descent-definition-descent-datum}{{023V}{source}{}{sourcepointer-8}{}}
\newlabel{descent-definition-descent-datum@cref}{{[originalreference][0][]023V}{[1][0][]source}{}{}{}}
\makeatother
\crefname{originalreference}{original reference}{original references}
\Crefname{originalreference}{Original reference}{Original references}
\makeatletter
\newlabel{more-morphisms-lemma-etale-separated-ind-quasi-affine}{{0AP9}{source}{}{sourcepointer-9}{}}
\newlabel{more-morphisms-lemma-etale-separated-ind-quasi-affine@cref}{{[originalreference][0][]0AP9}{[1][0][]source}{}{}{}}
\makeatother
\crefname{originalreference}{original reference}{original references}
\Crefname{originalreference}{Original reference}{Original references}
\makeatletter
\newlabel{morphisms-lemma-etale-locally-quasi-finite}{{03WS}{source}{}{sourcepointer-10}{}}
\newlabel{morphisms-lemma-etale-locally-quasi-finite@cref}{{[originalreference][0][]03WS}{[1][0][]source}{}{}{}}
\makeatother
\crefname{originalreference}{original reference}{original references}
\Crefname{originalreference}{Original reference}{Original references}
\makeatletter
\newlabel{more-groupoids-lemma-ind-quasi-affine}{{0APK}{source}{}{sourcepointer-11}{}}
\newlabel{more-groupoids-lemma-ind-quasi-affine@cref}{{[originalreference][0][]0APK}{[1][0][]source}{}{}{}}
\makeatother
\crefname{originalreference}{original reference}{original references}
\Crefname{originalreference}{Original reference}{Original references}
\makeatletter
\newlabel{descent-lemma-descending-property-surjective}{{02KV}{source}{}{sourcepointer-12}{}}
\newlabel{descent-lemma-descending-property-surjective@cref}{{[originalreference][0][]02KV}{[1][0][]source}{}{}{}}
\makeatother
\crefname{originalreference}{original reference}{original references}
\Crefname{originalreference}{Original reference}{Original references}
\makeatletter
\newlabel{descent-lemma-descending-property-etale}{{02VN}{source}{}{sourcepointer-13}{}}
\newlabel{descent-lemma-descending-property-etale@cref}{{[originalreference][0][]02VN}{[1][0][]source}{}{}{}}
\makeatother
\crefname{originalreference}{original reference}{original references}
\Crefname{originalreference}{Original reference}{Original references}
\makeatletter
\newlabel{descent-lemma-fpqc-universal-effective-epimorphisms}{{023Q}{source}{}{sourcepointer-14}{}}
\newlabel{descent-lemma-fpqc-universal-effective-epimorphisms@cref}{{[originalreference][0][]023Q}{[1][0][]source}{}{}{}}
\makeatother
\title{Algebraic spaces satisfy the sheaf property for the fpqc topology (Gabber)}
\author{The Stacks Project Authors}
\date{stacks-project:spaces-properties.tex:local-snapshot-20261012; GNU FDL 1.2 or later}
\begin{document}
\maketitle
\noindent\textit{Editorial note.} All delivered mathematical text is the author's own, copied verbatim from the pinned Stacks Project chapter file spaces-properties.tex; no statement or proof step was rewritten, reordered or removed.\par\medskip
\noindent\textit{Editorial note.} Span recomputation by label pairing (LANE\_BRIEF 10.2): in this chapter the author's proof environment follows the statement environment as a standalone block; the spec covers the complete environment pair (statement span, proof span) as recorded in the pool candidate.\par\medskip
\noindent\textit{Editorial note.} Occupation scan: combined claim\_key + statement\_anchor + generic-line scan over all registration files (active and withdrawn) found 13 S10-geom published cards for spaces-properties.tex (prop-occupied.json, 2026-10-14); those 13 labels are excluded from this batch by design. For the 77 free labels: 0 occupied. The registry mentions field (7557 files) and focused-seeds (645) both point to the same S10-geom cards (dual-source check), no other occupation.\par\medskip
\ifdefined\hypertarget\hypertarget{sourcepointer-1}{}\fi
\noindent\textit{Source reference.} Original reference 022H: sources/stacks-chow/tags.txt, source lines 2706--2706. Cross-chapter standard prerequisite cited by the author inside the counted span; it is used as a pointer and its own proof is not part of this item.\par\medskip
\ifdefined\hypertarget\hypertarget{sourcepointer-2}{}\fi
\noindent\textit{Source reference.} Original reference 02X4: sources/stacks-chow/tags.txt, source lines 3827--3827. Cross-chapter standard prerequisite cited by the author inside the counted span; it is used as a pointer and its own proof is not part of this item.\par\medskip
\ifdefined\hypertarget\hypertarget{sourcepointer-3}{}\fi
\noindent\textit{Source reference.} Original reference 02L4: sources/stacks-chow/tags.txt, source lines 3395--3395. Cross-chapter standard prerequisite cited by the author inside the counted span; it is used as a pointer and its own proof is not part of this item.\par\medskip
\ifdefined\hypertarget\hypertarget{sourcepointer-4}{}\fi
\noindent\textit{Source reference.} Original lemma 03BX: sources/stacks-chow/tags.txt, source lines 4364--4364. Cross-chapter standard prerequisite cited by the author inside the counted span; it is used as a pointer and its own proof is not part of this item.\par\medskip
\ifdefined\hypertarget\hypertarget{sourcepointer-5}{}\fi
\noindent\textit{Source reference.} Original lemma 03FX: sources/stacks-chow/tags.txt, source lines 4508--4508. Cross-chapter standard prerequisite cited by the author inside the counted span; it is used as a pointer and its own proof is not part of this item.\par\medskip
\ifdefined\hypertarget\hypertarget{sourcepointer-6}{}\fi
\noindent\textit{Source reference.} Original lemma 03FZ: sources/stacks-chow/tags.txt, source lines 4510--4510. Cross-chapter standard prerequisite cited by the author inside the counted span; it is used as a pointer and its own proof is not part of this item.\par\medskip
\ifdefined\hypertarget\hypertarget{sourcepointer-7}{}\fi
\noindent\textit{Source reference.} Original lemma 03H4: sources/stacks-chow/tags.txt, source lines 4553--4553. Cross-chapter standard prerequisite cited by the author inside the counted span; it is used as a pointer and its own proof is not part of this item.\par\medskip
\ifdefined\hypertarget\hypertarget{sourcepointer-8}{}\fi
\noindent\textit{Source reference.} Original reference 023V: sources/stacks-chow/tags.txt, source lines 2756--2756. Cross-chapter standard prerequisite cited by the author inside the counted span; it is used as a pointer and its own proof is not part of this item.\par\medskip
\ifdefined\hypertarget\hypertarget{sourcepointer-9}{}\fi
\noindent\textit{Source reference.} Original reference 0AP9: sources/stacks-chow/tags.txt, source lines 13467--13467. Cross-chapter standard prerequisite cited by the author inside the counted span; it is used as a pointer and its own proof is not part of this item.\par\medskip
\ifdefined\hypertarget\hypertarget{sourcepointer-10}{}\fi
\noindent\textit{Source reference.} Original reference 03WS: sources/stacks-chow/tags.txt, source lines 5123--5123. Cross-chapter standard prerequisite cited by the author inside the counted span; it is used as a pointer and its own proof is not part of this item.\par\medskip
\ifdefined\hypertarget\hypertarget{sourcepointer-11}{}\fi
\noindent\textit{Source reference.} Original reference 0APK: sources/stacks-chow/tags.txt, source lines 13478--13478. Cross-chapter standard prerequisite cited by the author inside the counted span; it is used as a pointer and its own proof is not part of this item.\par\medskip
\ifdefined\hypertarget\hypertarget{sourcepointer-12}{}\fi
\noindent\textit{Source reference.} Original reference 02KV: sources/stacks-chow/tags.txt, source lines 3386--3386. Cross-chapter standard prerequisite cited by the author inside the counted span; it is used as a pointer and its own proof is not part of this item.\par\medskip
\ifdefined\hypertarget\hypertarget{sourcepointer-13}{}\fi
\noindent\textit{Source reference.} Original reference 02VN: sources/stacks-chow/tags.txt, source lines 3774--3774. Cross-chapter standard prerequisite cited by the author inside the counted span; it is used as a pointer and its own proof is not part of this item.\par\medskip
\ifdefined\hypertarget\hypertarget{sourcepointer-14}{}\fi
\noindent\textit{Source reference.} Original reference 023Q: sources/stacks-chow/tags.txt, source lines 2751--2751. Cross-chapter standard prerequisite cited by the author inside the counted span; it is used as a pointer and its own proof is not part of this item.\par\medskip
\begin{proposition}[Gabber]
\label{proposition-sheaf-fpqc}
Let $S$ be a scheme. Let $X$ be an algebraic space over $S$. Then
$X$ satisfies the sheaf property for the fpqc topology.
\end{proposition}

\begin{proof}
Since $X$ is a sheaf for the Zariski topology it suffices to show
the following. Given a surjective flat morphism of affines
$f : T' \to T$ we have:
$X(T)$ is the equalizer of the two maps $X(T') \to X(T' \times_T T')$.
See Topologies, Lemma \ref{topologies-lemma-sheaf-property-fpqc}
(there is a little argument omitted here because the lemma cited
is formulated for functors defined on the category of all schemes).

\medskip\noindent
Let $a, b : T \to X$ be two morphisms such that $a \circ f = b \circ f$.
We have to show $a = b$. Consider the fibre product
$$
E = X \times_{\Delta_{X/S}, X \times_S X, (a, b)} T.
$$
By Spaces, Lemma \ref{spaces-lemma-properties-diagonal}
the morphism $\Delta_{X/S}$ is a representable monomorphism. Hence
$E \to T$ is a monomorphism of schemes. Our assumption that
$a \circ f = b \circ f$ implies that $T' \to T$ factors (uniquely) through $E$.
Consider the commutative diagram
$$
\xymatrix{
T' \times_T E \ar[r] \ar[d] & E \ar[d] \\
T' \ar[r] \ar@/^5ex/[u] \ar[ru] & T
}
$$
Since the projection $T' \times_T E \to T'$ is a monomorphism
with a section we conclude it is an isomorphism. Hence we conclude that
$E \to T$ is an isomorphism by
Descent, Lemma \ref{descent-lemma-descending-property-isomorphism}.
This means $a = b$ as desired.

\medskip\noindent
Next, let $c : T' \to X$ be a morphism such that the two compositions
$T' \times_T T' \to T' \to X$ are the same. We have to find a morphism
$a : T \to X$ whose composition with $T' \to T$ is $c$. Choose an
affine scheme $U$ and an \'etale morphism $U \to X$ such that the image
of $|U| \to |X|$ contains the image of $|c| : |T'| \to |X|$.
This is possible by Lemmas \ref{lemma-topology-points} and
\ref{lemma-cover-by-union-affines}, the fact that a finite disjoint union of
affines is affine, and the fact that $|T'|$ is quasi-compact
(small argument omitted). Since $U \to X$ is separated
(Lemma \ref{lemma-separated-cover}), we see that
$$
V = U \times_{X, c} T' \longrightarrow T'
$$
is a surjective, \'etale, separated morphism of schemes
(to see that it is surjective use Lemma \ref{lemma-points-cartesian}
and our choice of $U \to X$). The fact that
$c \circ \text{pr}_0 = c \circ \text{pr}_1$ means that we obtain a
descent datum on $V/T'/T$
(Descent, Definition \ref{descent-definition-descent-datum})
because
\begin{align*}
V \times_{T'} (T' \times_T T')
& =
U \times_{X, c \circ \text{pr}_0} (T' \times_T T') \\
& =
(T' \times_T T') \times_{c \circ \text{pr}_1, X} U \\
& =
(T' \times_T T') \times_{T'} V
\end{align*}
The morphism $V \to T'$ is ind-quasi-affine by
More on Morphisms, Lemma
\ref{more-morphisms-lemma-etale-separated-ind-quasi-affine}
(because \'etale morphisms are locally quasi-finite, see
Morphisms, Lemma \ref{morphisms-lemma-etale-locally-quasi-finite}).
By More on Groupoids, Lemma \ref{more-groupoids-lemma-ind-quasi-affine}
the descent datum is effective. Say $W \to T$ is a morphism
such that there is an isomorphism $\alpha : T' \times_T W \to V$
compatible with the given descent datum on $V$ and the canonical descent
datum on $T' \times_T W$. Then $W \to T$ is surjective and \'etale
(Descent, Lemmas \ref{descent-lemma-descending-property-surjective} and
\ref{descent-lemma-descending-property-etale}).
Consider the composition
$$
b' : T' \times_T W \longrightarrow V = U \times_{X, c} T' \longrightarrow U
$$
The two compositions
$b' \circ (\text{pr}_0, 1), 
b' \circ (\text{pr}_1, 1) :
(T' \times_T T') \times_T W \to T' \times_T W \to U$
agree by our choice of $\alpha$ and the corresponding property of $c$
(computation omitted). Hence $b'$ descends to a morphism $b : W \to U$ by
Descent, Lemma \ref{descent-lemma-fpqc-universal-effective-epimorphisms}.
The diagram
$$
\xymatrix{
T' \times_T W \ar[r] \ar[d] & W \ar[r]_b & U \ar[d] \\
T' \ar[rr]^c &  & X
}
$$
is commutative. What this means is that we have proved the existence
of $a$ \'etale locally on $T$, i.e., we have an $a' : W \to X$.
However, since we have proved uniqueness
in the first paragraph, we find that this \'etale local solution
satisfies the glueing condition, i.e., we have
$\text{pr}_0^*a' = \text{pr}_1^*a'$ as elements of $X(W \times_T W)$.
Since $X$ is an \'etale sheaf we find a unique $a \in X(T)$ restricting
to $a'$ on $W$.
\end{proof}

\end{document}
